% Written by projects/explain-top-move/paper/generate.mjs; do not edit by hand.
\begin{minipage}[t]{0.4\textwidth}\vspace{0pt}
\chessboard[setfen={rn1q2k1/ppp3p1/3b1p1p/3p1R1Q/3P3B/2P1r3/PP1N2PP/R5K1 w - - 0 17}, boardfontsize=13pt, pgfstyle=straightmove, arrow=to, linewidth=0.07em, color=red!75!black, markmoves={a1-f1}, pgfstyle=straightmove, arrow=to, linewidth=0.07em, color=green!50!black, markmoves={f5-f6}]
\end{minipage}\hfill
\begin{minipage}[t]{0.56\textwidth}\vspace{0pt}\small
\textbf{Played:} 17.~\symrook{}af1, White to move; the app's verdict: Blunder.\par
\textbf{Engine:} +2.5 with the best move, 0.0 after yours.\par
\textbf{Best line:} 17.~\symrook{}xf6 \symqueen{}e8 18.~\symqueen{}xd5+ \symking{}h8 19.~\symrook{}f7 \symknight{}c6\par
\textbf{UCI:} \uci{a1f1}; best line \uci{f5f6 d8e8 h5d5 g8h8 f6f7 b8c6}\par
\textbf{FEN:} \fen{rn1q2k1/ppp3p1/3b1p1p/3p1R1Q/3P3B/2P1r3/PP1N2PP/R5K1 w - - 0 17}
\end{minipage}

\smallskip\small
\begin{tabularx}{\textwidth}{>{\raggedright\arraybackslash}p{0.17\textwidth} L L}
\toprule
section & before & after \\
\midrule
What Rxf6 does instead & it attacks the bishop on d6 and the pawn on h6 \textperiodcentered{} 17.Rxf6 Qe8 18.Qxd5+ Kh8 19.Rf7 Nc6 20.Raf1 Ne7 & it attacks the bishop on d6 and the pawn on h6; if they take it with gxf6, Qg6+ Kf8 Qxh6+ Kf7, then Qxe3 takes their rook on e3, and you are two pawns' worth of material up; if they take it with Qxf6, Bxf6 wins their queen for your rook \textperiodcentered{} 17.Rxf6 Qe8 18.Qxd5+ Kh8 19.Rf7 Nc6 20.Raf1 Ne7 \\
Retry answer & The move was Rxf6 (the green arrow). It attacks the bishop on d6 and the pawn on h6. & The move was Rxf6 (the green arrow). It attacks the bishop on d6 and the pawn on h6. If they take it with gxf6, Qg6+ Kf8 Qxh6+ Kf7, then Qxe3 takes their rook on e3, and you are two pawns' worth of material up. If they take it with Qxf6, Bxf6 wins their queen for your rook. \\
\bottomrule
\end{tabularx}

\textbf{The judge (run 1, where the new text was X):} ``X shows why the rook sacrifice works (gxf6 Qg6+ Kf8 Qxh6+ then Qxe3 wins the rook back, and Qxf6 Bxf6 wins the queen), while Y leaves the player wondering why Rxf6 does not just lose a rook to gxf6.''
